For \(z\in\mathcal S_k\),
\[
 \frac{\gamma_k(z)}{\beta_k(z)}
 =\frac{M_k^{-1}}{\pi_k^{m_k}(1-\pi_k)^{n-m_k}}=\frac1{B_{kk}}. \tag{22}
\]
Substitution in the Horvitz--Thompson formula, with \(\eta_k=p_kB_{kk}>0\), proves the displayed estimator identity.  Its iid summand has mean \(U_k\), because conditional assignment is uniform on \(\mathcal S_k\), and hence its exact influence coordinate is
\(\phi_k^{HOS}=\eta_k^{-1}\mathbf1\{L=k,Z\in\mathcal S_k\}W(Z)-U_k\).  Furthermore
\[
 \mathbb E\{W(Z)^2\mid L=k,Z\in\mathcal S_k\}
 =\tau_k^2+M_k^{-1}\sum_{z\in\mathcal S_k}\mu_z^2
 =\tau_k^2+b_k^2+U_k^2.
\]
Therefore
\[
 \operatorname{Var}(\phi_k^{HOS})
 =\frac{\tau_k^2+b_k^2}{\eta_k}+U_k^2(\eta_k^{-1}-1). \tag{23}
\]
In the nominal-label target-hit subexperiment, each \(z\in\mathcal S_k\) occurs with probability \(\eta_k/M_k\).  The tangent-space calculation (10)--(12), with \(q_k\) replaced by \(\eta_k\), proves that \(\phi_k^{nom}\) is canonical and has variance \(\tau_k^2/\eta_k\).  Since \(b_k^2,U_k^2\geq0\) and \(0<\eta_k\leq1\), (23) is at least this bound.  If \(q_k>\eta_k\) and \(\tau_k^2>0\), positivity permits division and gives
\[
 \tau_k^2/q_k<\tau_k^2/\eta_k\leq\operatorname{Var}(\phi_k^{HOS}).
\]
Finally \(b_k^2>0\) makes the cell-heterogeneity term \(b_k^2/\eta_k\) strictly positive, independently of the nonnegative participation term.